\section{Related Work}\label{sec:related}

\subsection{Recurrent networks for cryptocurrency forecasting}

Recurrent neural networks keep a hidden state that carries information from one time step to the next, which makes them a natural fit for price series. Plain RNNs struggle with long sequences because gradients shrink as they are passed back through many steps, so the network forgets early events. The LSTM \cite{hochreiter1997} adds a cell state and three gates that control what is kept and what is dropped, and the GRU \cite{cho2014} simplifies this to two gates and a single state, which makes it smaller and faster to train. A bidirectional network \cite{schuster1997} runs one recurrent layer forwards and another backwards over the input window and combines the two.

Seabe et al. \cite{seabe2023} compared the LSTM, GRU and Bi-LSTM on Bitcoin, Ethereum and Litecoin and found that the Bi-LSTM had the lowest error, with a Bitcoin MAPE of 0.0356 against 0.0394 for the LSTM. Hansun et al. \cite{hansun2022} gave the same three networks several inputs per coin (the open, high, low, close and volume) for five cryptocurrencies, and found that the Bi-LSTM and GRU performed about equally well. Both studies judge models by forecast error. A forecast error on its own is hard to read unless it is set next to a naive forecast that simply repeats today's price, because daily prices move very little relative to their level, so even a model with no skill can post a small MAPE. Forecast error also says nothing about whether the forecasts would make or lose money once costs and leverage are included, which is the question this project asks.

\subsection{Attention and multivariate inputs}

Attention lets a model weight each step of its input by how relevant it is to the current output, instead of relying only on the final hidden state. It was introduced for translation \cite{bahdanau2015} and became the core of the Transformer \cite{vaswani2017}. Adding an attention layer on top of a recurrent network is a common middle ground, since it keeps the recurrent memory but lets the model look back directly at, for example, the day of a crash.

There are good reasons to give a Bitcoin model more than Bitcoin's own history. Bitcoin leads the rest of the cryptocurrency market \cite{corbet2018}, and shocks spread quickly between coins \cite{huynh2020}. Bitcoin's link with equities changes over time and it behaves more like a diversifier than a hedge \cite{bouri2017}, while gold acts as a safe haven that capital moves into during crypto stress \cite{huynh2020}. Chen et al. \cite{chen2021ml} found that adding economic indicators and on-chain data such as hash rate and transaction volume improved Bitcoin forecasts, and Liu and Tsyvinski \cite{liu2021} found that network adoption measures and momentum predict cryptocurrency returns at weekly horizons. The halving, which cuts the reward for mining roughly every four years, is also widely believed to move prices and market sentiment \cite{subawa2025}, and the market has gone through very different phases between halvings. That is a reason to train across several years rather than one.

\subsection{Volatility and jump risk}

Returns are hard to forecast, but their size is not. Large moves tend to be followed by large moves, a pattern known as volatility clustering \cite{cont2001}, and the GARCH(1,1) model \cite{bollerslev1986} captures it with a simple recursion. In a comparison of 330 volatility models, Hansen and Lunde \cite{hansen2005} found nothing that beat it on exchange-rate data, although models that allow for asymmetric responses did better on stock returns. Daily volatility cannot be observed directly, so forecasts are usually scored against a proxy such as the absolute or squared return. Such proxies are very noisy, which makes every forecast look worse than it is. Patton \cite{patton2011} showed that only some ways of scoring forecasts against a noisy proxy still rank them correctly, so the choice of target and loss needs care.

Bitcoin also jumps. Its price does not always move smoothly through intermediate levels, and jump-diffusion models add a separate process for these sudden moves \cite{xu2019}. Chen and Huang \cite{chen2021jump} detected jumps in Bitcoin prices with the Lee and Mykland test, which compares each return with a rolling estimate of normal variation, and found that jumps in both returns and volatility are significant in Bitcoin's price history. This project uses a simpler threshold rule in the same spirit (\cref{sec:data}). Azizi et al. \cite{azizi2023} used a neural network to estimate Bitcoin's volatility from historical data, which is the same task the Beta Agent takes on.

\subsection{Risk management for leveraged trading}

Scaling a position down when volatility is high and up when it is low is an old idea. Moreira and Muir \cite{moreira2017} showed that portfolios which take less risk when volatility is high earn higher Sharpe ratios, for the stock market, several equity factors and the currency carry trade, because changes in volatility are not matched by changes in expected returns. Stops that scale with recent volatility are the same idea applied to exits. Results are normally judged by maximum drawdown and the Sharpe ratio \cite{sharpe1994}, with the Sortino ratio \cite{sortino1994} as a variant that only penalises downside moves.

Backtests are easy to fool. Trying many settings on the same data and keeping the best one produces results that do not survive on new data \cite{bailey2014}, which is why the tuning, testing and forward-testing periods in this project are kept apart. Leverage on crypto CFDs is also regulated. Since 2018 retail clients in the EU have been limited to 2:1 \cite{esma2018}, and since 2021 the UK has banned the sale of crypto derivatives to retail clients altogether \cite{fca2020}.

Put together, this literature has strong forecasting models and a well understood toolkit for risk management, but the two are rarely joined. Studies of LSTM-type models for cryptocurrency mostly stop at forecast error, and studies of volatility management are mostly about unleveraged portfolios. Few evaluate a learned volatility forecast that drives sizing, stops and leverage, end to end, under the conditions a retail CFD trader faces. That is the gap this project addresses.
